\documentclass[10pt,a4paper]{article}
\usepackage{fontspec,amsmath,amssymb,graphicx}
\usepackage[margin=18mm]{geometry}
\setmainfont{Times New Roman}
\setlength{\parindent}{0pt}
\setlength{\parskip}{5pt}
\allowdisplaybreaks
\sloppy
\emergencystretch=2em
\begin{document}
\section*{VD\_2\_2\_09 — Lời giải bổ sung}
\subsection*{Kiến thức và trình tự giải}
\textbf{Mục tiêu.} Giải phương trình hyperbolic sau khi đã tìm hai tọa độ đặc trưng độc lập.

\textbf{Bước 1 — Phân loại, tìm đặc trưng và thay đạo hàm.} Làm đủ các bước của phép đưa về chính tắc. Viết cả phép đổi ngược để biểu diễn vế phải $\widetilde f(\xi,\eta)=f(x(\xi,\eta),y(\xi,\eta))$.

\textbf{Bước 2 — Tích phân phương trình chính tắc.} Với $v_{\xi\eta}=g$, giữ biến kia cố định khi tích phân:


\[
v=F(\xi)+G(\eta)+\int_{\xi_0}^{\xi}\int_{\eta_0}^{\eta}g(a,b)\,db\,da.
\]


Nếu tìm được một nghiệm riêng $P(x,y)$ dễ hơn, viết $u=P+F(\xi)+G(\eta)$. Với $v_{\xi\eta}+\alpha v_\eta=0$, đặt $w=v_\eta$, giải $w_\xi+\alpha w=0$ rồi tích phân theo $\eta$; thu được $v=F(\xi)+e^{-\alpha\xi}G(\eta)$.

\textbf{Bước 3 — Đổi cả dữ liệu.} Thế tọa độ đường dữ liệu vào $\xi,\eta$, dùng quy tắc dây chuyền cho đạo hàm được cho. Trên đường không đặc trưng, hai dữ liệu thường xác định hai hàm $F,G$. Trên đường đặc trưng, có biến không thay đổi, nên dữ liệu chỉ xác định giá trị hoặc đạo hàm của một hàm tại một điểm; phải tìm điều kiện tương thích và nêu phần tùy ý còn lại.

\textbf{Bước 4 — Lập đẳng thức hình bình hành.} Với $v=F(\xi)+G(\eta)$, bốn đỉnh theo hai họ đặc trưng thỏa $v_{11}+v_{22}=v_{12}+v_{21}$. Nếu có số hạng mũ, trước hết nhân trọng số để được tổng hai hàm một biến. Với vế phải $g$, hiệu bốn đỉnh bằng tích phân kép của $g$ trên hình chữ nhật trong tọa độ mới.

\textbf{Bước 5 — Xác định miền được dữ liệu quyết định.} Hai đặc trưng đi qua điểm đang xét phải gặp đoạn dữ liệu phù hợp. Với phương trình sóng $v_{tt}=v_{zz}$ trên $z>0$, dữ liệu lúc $t=0$ xác định duy nhất miền $z>t$; ở $0<z<t$, cần thêm dữ liệu biên $z=0$. Với đặc trưng cong, dùng miền trong các tọa độ mới rồi đổi lại.

\textbf{Bước 6 — Kiểm tra.} Thế nghiệm vào cả PDE và từng điều kiện. Nếu đề yêu cầu hai nghiệm, chỉ rõ hai lựa chọn khác nhau cho hàm tùy ý và chứng minh phần chênh lệch thỏa phương trình thuần nhất, đồng thời triệt tiêu trên đường dữ liệu.

\textbf{Biến đổi riêng của bài.} Dùng tọa độ cực (hoặc cực sau khi co giãn trục). Khi thu được $r v_{r\theta}-v_\theta=0$, đặt $w=v/r$ thì $w_{r\theta}=0$. Vì vậy $v=rF(\theta)+G(r)$; phải chia nghiệm cho $r$ trước khi lập đẳng thức bốn đỉnh.
\subsection*{Lời giải áp dụng cho bài}
\textbf{(i) Phân loại và tọa độ cực.} Với $A=xy$, $2B=y^2-x^2$, $C=-xy$,
\[\delta=\frac{(y^2-x^2)^2}{4}+x^2y^2=\frac{(x^2+y^2)^2}{4}>0\quad(y>0).\]
Đặc trưng phân tích thành $(xdy-ydx)(ydy+xdx)=0$: các tia $x/y=\mathrm{const}$ và nửa đường tròn $x^2+y^2=\mathrm{const}$. Chọn $r=\sqrt{x^2+y^2}>0$, $\theta\in(0,\pi)$, $x=r\cos\theta$, $y=r\sin\theta$. Phép đổi ngược có Jacobian $r\ne0$.

\textbf{Tính toán phần chính.} Thế ba công thức đạo hàm cực ở Ví dụ 5; các hạng $v_{rr},v_{\theta\theta},v_r$ triệt tiêu, còn
\[xyu_{xx}+(y^2-x^2)u_{xy}-xyu_{yy}=-r v_{r\theta}+v_\theta.\]
Dạng chính tắc là $r v_{r\theta}-v_\theta=0$.

\textbf{(ii) Nghiệm tổng quát.} Đặt $w=v/r$. Vì $v_{r\theta}=w_\theta+r w_{r\theta}$, phương trình cho $r^2w_{r\theta}=0$. Tích phân hai lần: $w=F(\theta)+K(r)$, tức
\[\boxed{u=rF(\theta)+G(r).}\]

\textbf{(iii) Dữ liệu $y=1$.} Từ $u(x,1)=x^2$ có $u_x(x,1)=2x$; kết hợp $u_y(x,1)=4$, quy tắc dây chuyền cho
\[v_\theta=-y u_x+xu_y=2x\quad(y=1).\]
Trên đường này $r=1/\sin\theta$, $x=\cot\theta$, nên $rF'(\theta)=2\cot\theta$ và $F'=2\cos\theta$. Lấy $F=2\sin\theta$ (hằng số được hấp thụ vào $G$), dữ liệu giá trị cho $G(r)=r^2-3$, $r\ge1$. Nghiệm cụ thể
\[\boxed{u=x^2+y^2+2y-3}\]
thỏa PDE vì $u_{xx}=u_{yy}=2$, $u_{xy}=0$, và thỏa cả hai dữ liệu. Công thức đa thức xác định trên toàn $y>0$, nhưng dữ liệu chỉ quyết định duy nhất miền $x^2+y^2>1$, $y>0$ (và giá trị giới hạn tại $r=1$). Ở $r<1$ có thể đổi $G$ bằng một hàm $C^2$ triệt tiêu cùng hai đạo hàm tại $r=1$.

\textbf{(iv) Dữ liệu trên nửa đường tròn đơn vị.} Pháp tuyến ngoài là $\nu=(x,y)$, nên $\partial_\nu u=v_r$ tại $r=1$. Dữ liệu cho
\[F(\theta)+G(1)=\varphi(\cos\theta),\qquad F(\theta)+G'(1)=\psi(\cos\theta).\]
Điều kiện cần và đủ là $\psi(x)-\varphi(x)=C$ hằng trên $(-1,1)$. Toàn bộ nghiệm có thể viết
\[u=r\varphi(x/r)+H(r),\qquad H(1)=0,\quad H'(1)=C.\]
Ví dụ $H=C(r-1)$ hoặc $H=C(r-1)+(r-1)^2$; vì thế dữ liệu trên một đặc trưng không xác định duy nhất. Giả sử $\varphi$ và $H$ đủ trơn trên miền xét.

\textbf{(v) Hai dữ liệu giá trị và ghép nghiệm.} Trên $r=1$ có $F(\theta)+G(1)=1$ cho mọi $\theta$, nên $F$ hằng; hấp thụ nó vào $G$ được $u=H(r)$. Dữ liệu $u(x,0)=x$, $1\le x\le2$ cho $H(r)=r$ ở $1\le r\le2$. Vậy
\[\boxed{u=\sqrt{x^2+y^2}\quad\text{trên }1<x^2+y^2<4,\ y>0.}\]
Đây là miền duy nhất do hai dữ liệu quyết định. Với $0<r\le1$, ghép $u=h(r)$: các điều kiện $h(1)=1$, $h'(1)=1$, $h''(1)=0$ làm giá trị và các đạo hàm đến cấp hai khớp tại $r=1$. Mọi hàm chỉ phụ thuộc $r$ đều thỏa PDE, nên phép ghép cho nghiệm $C^2$ trên $0<r<2$, $y>0$, liên tục đến phần biên đã cho. Ngoài $r=2$, có thể tiếp tục bằng $u=r$ hoặc một phần ghép đủ trơn khác.

\textbf{Chi tiết gom năm hệ số trong tọa độ cực.} Viết $c=\cos\theta,s=\sin\theta$, $xy=r^2cs$, $y^2-x^2=r^2(s^2-c^2)$. Dùng ba đạo hàm cực, hệ số lần lượt là
\[\begin{aligned}
v_{rr}:&\quad r^2[cs\,c^2+(s^2-c^2)cs-cs\,s^2]=0,\\
v_{r\theta}:&\quad r[-2c^2s^2-(c^2-s^2)^2-2c^2s^2]=-r,\\
v_{\theta\theta}:&\quad cs\,s^2-(s^2-c^2)cs-cs\,c^2=0,\\
v_r:&\quad r[cs\,s^2-(s^2-c^2)cs-cs\,c^2]=0,\\
v_\theta:&\quad 2c^2s^2+(c^2-s^2)^2+2c^2s^2=1.
\end{aligned}\]
Đã dùng $4c^2s^2+(c^2-s^2)^2=(c^2+s^2)^2=1$. Vì vậy toán tử đúng bằng $-rv_{r\theta}+v_\theta$, kể cả tại các hướng có $x=0$ trong miền $y>0$.\end{document}